\documentclass[11pt]{article}
\usepackage[hyperref]{acl}
\usepackage{times}
\usepackage{latexsym}
\usepackage{amsmath}
\usepackage{amssymb}
\usepackage{graphicx}
\usepackage{booktabs}
\usepackage{multirow}
\usepackage{xcolor}
\usepackage{colortbl}
\usepackage{tabularx}
\usepackage{enumitem}

\title{Measuring Policy Ambiguity's Causal Impact on Agent Safety:\\ A Matched-Pair Framework with Judge Calibration}

\author{Anonymous}

\begin{document}
\maketitle

\begin{abstract}
LLM-based agents increasingly operate under natural-language policy specifications that inevitably contain ambiguity.
We present the first matched-pair causal study measuring how six types of policy ambiguity---spanning specification-level (authorization scope, incompleteness, conditional precedence) and linguistic-level (scopal, lexical, coreferential) categories---affect safety violations across five architecturally diverse LLM agents.
Using 300 ambiguous/unambiguous clause pairs injected into the $\tau^2$-bench environment, we collect 2{,}997 agent episodes and apply a two-stage evaluation pipeline with LLM-as-judge assessment followed by human-anchored calibration.
Our judge calibration study (100 human-annotated episodes) reveals that LLM judges systematically overestimate violation rates (F1\,=\,0.64, $\kappa$\,=\,0.56 vs.\ humans) despite excellent inter-judge agreement ($\kappa$\,=\,0.91) and test-retest reliability (ICC\,=\,0.919).
After correction, ambiguous policies still cause a substantial +19.5 percentage point increase in violations (OR\,=\,6.2, $p$\,$<$\,0.0001), confirmed as ambiguity-specific by perturbation controls ($\Delta_{\text{perturb}}$\,=\,$-$0.67pp).
We find that ambiguity types differ in severity ($\chi^2$\,=\,16.4, $p$\,=\,0.006) and that models exhibit distinct vulnerability profiles (model$\times$type interaction, $p$\,=\,0.028), with authorization-scope vulnerability ranging from 6\% (Claude) to 38\% (DeepSeek) after correction.
Qualitative analysis of 614 violation traces reveals pervasive \emph{ambiguity blindness} (4.6\% awareness rate, $\kappa$\,=\,0.71 vs.\ human anchor).
Our matched-pair methodology and judge calibration pipeline constitute reusable contributions for agent safety evaluation.
\end{abstract}

%------------------------------------------------------------
\section{Introduction}
%------------------------------------------------------------

Consider a customer-service agent operating under a return policy that states \emph{``items may be returned within 30 days.''}
The scopal ambiguity of ``items''---does it include customized products?---leads the agent to grant a full refund on a \$200 personalized order.
When the policy is disambiguated to ``standard items only,'' the same agent correctly refuses.
A single word of ambiguity causes an unauthorized financial action.

As LLM-based agents are deployed with natural-language policy specifications~\cite{yao2023react,schick2024toolformer}, the safety implications of policy ambiguity become urgent.
Prior work classifies ambiguity types in requirements~\cite{berry2003ambiguity,massey2014identifying,li2024drip} and evaluates agent safety under adversarial or missing policies~\cite{ruan2024toolemu,xie2024phantompolicy}, but no study causally isolates ambiguity's effect on safety violations at type-level granularity across diverse model families.

We address this gap with three contributions:

\paragraph{(1) Matched-pair causal framework.}
We design 300 ambiguous/unambiguous clause pairs spanning six ambiguity types organized in a two-layer taxonomy (specification vs.\ linguistic), inject them into $\tau^2$-bench~\cite{tau2bench2024} scenarios, and collect 2{,}997 episodes across five models (GPT-5.4, GPT-4.1, Claude Sonnet 4.6, Qwen3-235B, DeepSeek-V3).
Perturbation controls confirm the effect is ambiguity-specific ($\Delta_{\text{perturb}}$\,=\,$-$0.67pp vs.\ corrected ambiguity $\Delta$\,=\,+19.5pp).

\paragraph{(2) Judge calibration pipeline.}
We discover---and correct for---systematic LLM-judge inflation through a 100-episode human-annotated calibration study.
Despite excellent inter-judge reliability ($\kappa$\,=\,0.91, ICC\,=\,0.919), judges achieve only F1\,=\,0.64 against human ground truth ($\kappa$\,=\,0.56), revealing a reliability--accuracy dissociation with broad implications for LLM-as-judge evaluation pipelines.
After correction, the headline effect decreases from +34.0pp to +19.5pp (OR from 15.1 to 6.2) but remains highly significant ($p$\,$<$\,0.0001).

\paragraph{(3) Corrected differential analysis.}
After judge correction, we find that ambiguity types differ significantly in severity (corrected $\chi^2$\,=\,16.4, $p$\,=\,0.006, $V$\,=\,0.112), with incompleteness most dangerous (38\%) and coreferential least (16\%).
A model$\times$type interaction ($p$\,=\,0.028) reveals that authorization-scope vulnerability varies 32pp across models (6\%--38\%), while incompleteness remains universally dangerous (36\%--48\%).
Qualitative analysis of 614 violation traces identifies pervasive \emph{ambiguity blindness} (4.6\% awareness, human-validated $\kappa$\,=\,0.71).

%------------------------------------------------------------
\section{Related Work}
%------------------------------------------------------------

\paragraph{Ambiguity in requirements engineering.}
Berry and Kamsties~\cite{berry2003ambiguity} distinguish \emph{nocuous} ambiguity (divergent interpretations causing harm) from \emph{innocuous} ambiguity (convergent interpretations).
Massey et al.~\cite{massey2014identifying} identify six types of ambiguity in privacy policies.
Chantree et al.~\cite{chantree2006identifying} develop the first computational nocuous ambiguity detector.
Ferrari et al.~\cite{ferrari2017nlp} extend NLP-based detection across domains.
Our work operationalizes the nocuity concept for LLM agents: we measure whether ambiguity causes safety violations (nocuous) or is benignly resolved (innocuous).

\paragraph{Ambiguity in NLP.}
Li et al.~\cite{li2024drip} propose DRIP-R with 5 ambiguity types and show LLMs disagree on interpretations, but do not measure downstream safety impact.
Poesio and Artstein~\cite{poesio2005reliability} study annotation of anaphoric ambiguity.
Li et al.~\cite{li2024gardenpath} examine LLM comprehension of garden-path sentences.
Wasow et al.~\cite{wasow2005puzzle} address the ambiguity/vagueness boundary.
Our taxonomy draws from both NLP and RE traditions, with explicit cross-walks to prior classifications.

\paragraph{Agent safety evaluation.}
$\tau^2$-bench~\cite{tau2bench2024} evaluates tool-using agents on customer-service tasks.
ToolEmu~\cite{ruan2024toolemu} and R-Judge~\cite{yuan2024rjudge} assess agent safety under adversarial conditions.
PhantomPolicy~\cite{xie2024phantompolicy} studies missing/contradictory policies.
AgentPex~\cite{agentpex2024} examines procedural violations.
None of these isolates ambiguity type as a causal variable with matched-pair controls.

\paragraph{LLM-as-judge reliability.}
LLM judges are widely used for evaluation~\cite{zheng2024judging} but face known biases~\cite{wang2024large}.
Our calibration study provides the first systematic assessment of LLM-as-judge accuracy specifically for agent safety violation detection, revealing a reliability--accuracy dissociation not previously documented.

%------------------------------------------------------------
\section{Ambiguity Taxonomy}
\label{sec:taxonomy}
%------------------------------------------------------------

We organize policy ambiguity into two layers with six types, drawing on Li et al.~\cite{li2024drip} and Massey et al.~\cite{massey2014identifying}.

\paragraph{Specification layer} involves ambiguity in the policy's intended meaning, resolvable only by the policy author:
\begin{itemize}[nosep,leftmargin=*]
  \item \textbf{Authorization scope}: vague permission/role boundaries (e.g., ``authorized personnel'' without function specification).
  \item \textbf{Incompleteness}: missing conditions, exceptions, or edge cases.
  \item \textbf{Conditional precedence}: conflicting rules without priority ordering.
\end{itemize}

\paragraph{Linguistic layer} involves ambiguity arising from natural-language structure, partially resolvable via context:
\begin{itemize}[nosep,leftmargin=*]
  \item \textbf{Scopal}: quantifier, negation, or modal scope interactions.
  \item \textbf{Lexical}: polysemy or homonymy in policy terms.
  \item \textbf{Coreferential}: unclear antecedent resolution for pronouns or relative clauses.
\end{itemize}

The two-layer distinction reflects a \emph{disambiguation channel} asymmetry: linguistic ambiguities can sometimes be resolved through tool responses, world knowledge, or contextual inference, while specification ambiguities have no external resolution channel, forcing agents to guess.

%------------------------------------------------------------
\section{Method}
\label{sec:method}
%------------------------------------------------------------

\subsection{Clause Generation}
\label{sec:clause_gen}

We generate 300 matched clause pairs (50 per type) where each pair consists of an ambiguous clause and a minimally disambiguated variant.
Table~\ref{tab:clause_stats} summarizes per-type statistics.
Each pair was generated with a mix of 25 explicitly labeled and 25 naturalistic formulations.
Quality was assessed via automated screening with human spot-checks; the overall pass rate was 92\%.
Authorization scope (88\%) and coreferential (86\%) had the lowest pass rates, requiring 2--3 revision rounds due to inherent difficulty in cleanly separating boundary ambiguity from missing-permission ambiguity and referential from scopal overlap, respectively.

\begin{table}[t]
\centering
\small
\begin{tabular}{lcccc}
\toprule
\textbf{Type} & \textbf{Layer} & \textbf{Pass\%} & \textbf{$\bar{t}_\text{amb}$} & \textbf{$\Delta t$} \\
\midrule
Auth.\ scope & Spec & 88 & 28.3 & 3.9 \\
Incompleteness & Spec & 94 & 26.8 & 3.8 \\
Cond.\ preced. & Spec & 96 & 24.1 & 1.6 \\
Scopal & Ling & 96 & 22.6 & 1.4 \\
Lexical & Ling & 92 & 21.9 & 1.3 \\
Coreferential & Ling & 86 & 25.2 & 2.3 \\
\bottomrule
\end{tabular}
\caption{Per-type clause statistics. $\bar{t}_\text{amb}$: mean tokens in ambiguous clause. $\Delta t$: token delta (disambiguation length). Specification-layer clauses require more tokens to disambiguate (mean 33.6 vs.\ 11.9 for linguistic).}
\label{tab:clause_stats}
\end{table}

\subsection{Realism Validation}

To verify ecological validity, we compare linguistic properties of our synthetic clauses against natural ambiguities from DRIP-R~\cite{li2024drip}.
Table~\ref{tab:realism} shows Kolmogorov-Smirnov tests across four properties; none reaches significance at $p$\,$<$\,0.05, with parse depth marginal ($p$\,=\,0.052).
Synthetic clauses are slightly shorter (24.5 vs.\ 26.8 tokens mean) and shallower, consistent with the matched-pair design favoring minimal variation.

\begin{table}[t]
\centering
\small
\begin{tabular}{lcccc}
\toprule
\textbf{Property} & \textbf{Synth.} & \textbf{DRIP-R} & \textbf{KS} & \textbf{$p$} \\
\midrule
Sentence length & 24.5 & 26.8 & 0.14 & 0.19 \\
Parse depth & 3.2 & 3.9 & 0.21 & 0.052 \\
Flesch-Kincaid & 12.1 & 11.5 & 0.11 & 0.31 \\
Type-token ratio & 0.71 & 0.69 & 0.09 & 0.42 \\
\bottomrule
\end{tabular}
\caption{Realism comparison: synthetic clauses vs.\ DRIP-R natural ambiguities. No significant distributional differences at $p$\,$<$\,0.05.}
\label{tab:realism}
\end{table}

\subsection{Experimental Setup}

Clause pairs are injected into $\tau^2$-bench~\cite{tau2bench2024} customer-service scenarios spanning retail and airline domains.
Each of 300 clause pairs $\times$ 2 conditions (ambiguous/unambiguous) is evaluated on five models:

\begin{itemize}[nosep,leftmargin=*]
  \item \textbf{Proprietary}: GPT-5.4 (frontier), GPT-4.1 (mid-tier), Claude Sonnet 4.6
  \item \textbf{Open-source}: Qwen3-235B (MoE), DeepSeek-V3 (dense)
\end{itemize}

This yields 1{,}800 API-model episodes and 1{,}197 open-source episodes (3 errors in Qwen3), totaling 2{,}997.
Prior to the full study, a mini-pilot (20 episodes: 5 clauses $\times$ 2 types $\times$ 2 conditions on GPT-5.4) confirmed basic feasibility for the two new types (authorization scope, coreferential).

\subsection{Evaluation Pipeline}
\label{sec:eval_pipeline}

\paragraph{Stage 1: LLM-as-judge.}
Each episode is evaluated by a cross-model LLM judge (GPT-5.4 judges Claude episodes; GPT-4.1 judges the rest) for binary violation detection.
Inter-judge agreement is high ($\kappa$\,=\,0.91).

\paragraph{Stage 2: Human calibration.}
We stratify 100 episodes by type$\times$model$\times$condition and obtain binary violation labels from two human annotators (inter-annotator $\kappa$\,=\,0.79).
This reveals systematic judge inflation (\S\ref{sec:judge_cal}) and enables corrected estimates for all downstream claims.

\subsection{Perturbation Control}
\label{sec:perturbation}

To confirm ambiguity-specific causation, we construct 30 non-ambiguity perturbation pairs (synonym substitutions, clause reordering) and run 300 episodes.
The perturbation effect is $\Delta_{\text{perturb}}$\,=\,$-$0.67pp (95\% CI: [$-$8.0, 6.67], McNemar $p$\,=\,1.0), confirming that surface-level prompt changes produce negligible safety impact (Table~\ref{tab:perturbation}).

\begin{table}[t]
\centering
\small
\begin{tabular}{lccc}
\toprule
\textbf{Metric} & \textbf{Value} & \textbf{Threshold} & \textbf{Verdict} \\
\midrule
$\Delta_\text{perturb}$ & $-$0.67pp & $<$5pp & PASS \\
CI contains 0 & [$-$8.0, 6.67] & required & PASS \\
Ratio to ambiguity & 0.02 & $<$0.15 & PASS \\
\bottomrule
\end{tabular}
\caption{Perturbation control verdict. The near-zero perturbation effect ($\Delta$\,=\,$-$0.67pp) vs.\ the corrected ambiguity effect (+19.5pp) confirms ambiguity-specific causation.}
\label{tab:perturbation}
\end{table}

\subsection{Measurement Reliability}

Test-retest reliability is assessed via 30 clauses $\times$ 3 replicates $\times$ 2 conditions $\times$ 2 API models (360 episodes).
Overall ICC\,=\,0.919 [0.892, 0.940] (excellent); binary ICC\,=\,0.847.
Per-type ICCs range from 0.807 (incompleteness) to 0.989 (scopal).
As discussed in \S\ref{sec:judge_cal}, high ICC indicates \emph{reproducible} but not necessarily \emph{accurate} measurement---the judge consistently produces the same (inflated) assessment across replicates.

%------------------------------------------------------------
\section{Judge Calibration}
\label{sec:judge_cal}
%------------------------------------------------------------

The judge calibration study is a central methodological contribution.
We annotate 100 episodes ($\sim$17 per model, balanced by condition and type) with two human annotators (inter-annotator $\kappa$\,=\,0.79, 87\% agreement).

\subsection{Reliability--Accuracy Dissociation}

Table~\ref{tab:dissociation} reveals the core finding: LLM judges achieve excellent agreement with each other ($\kappa$\,=\,0.91) and excellent test-retest reliability (ICC\,=\,0.919), but only moderate agreement with human ground truth (F1\,=\,0.64, $\kappa$\,=\,0.56).
The 0.35 gap between LLM-LLM $\kappa$ (0.91) and LLM-human $\kappa$ (0.56) quantifies a shared systematic bias.

\begin{table}[t]
\centering
\small
\begin{tabular}{lccl}
\toprule
\textbf{Metric} & \textbf{Value} & \textbf{Source} & \textbf{Interpretation} \\
\midrule
LLM-LLM $\kappa$ & 0.91 & Cross-judge & Judges agree \\
Human-human $\kappa$ & 0.79 & Calibration & Task is clear \\
ICC test-retest & 0.919 & Reliability & Reproducible \\
LLM-human $\kappa$ & 0.56 & Calibration & \textbf{Accuracy gap} \\
LLM-human F1 & 0.64 & Calibration & Moderate accuracy \\
\bottomrule
\end{tabular}
\caption{Reliability vs.\ accuracy dissociation. Judges are reproducible and self-consistent but systematically deviate from human ground truth.}
\label{tab:dissociation}
\end{table}

\subsection{Inflation Mechanism}

The inflation is condition-asymmetric: the judge achieves F1\,=\,0.73 and FPR\,=\,0.05 on unambiguous episodes but only F1\,=\,0.61 and FPR\,=\,0.19 on ambiguous ones (Table~\ref{tab:judge_condition}).
The net delta inflation is 10.4pp (13.2pp ambiguous $-$ 2.8pp unambiguous).

The root cause is that the LLM judge applies a strict ground-truth standard: when an agent acts on a \emph{plausible} interpretation of an ambiguous clause, human annotators judge $\sim$40\% of such cases as ``reasonable under ambiguity'' (non-violation), but the judge---which has access to the unambiguous ground truth---labels any deviation as a violation.
The judge lacks the concept of ``reasonable ambiguity resolution.''

\begin{table}[t]
\centering
\small
\begin{tabular}{lcccc}
\toprule
\textbf{Condition} & \textbf{F1} & \textbf{FPR} & \textbf{Inflation} & \textbf{Mechanism} \\
\midrule
Ambiguous & 0.61 & 0.19 & +13.2pp & No ``reasonable resolution'' \\
Unambiguous & 0.73 & 0.05 & +2.8pp & Clear-cut cases \\
\bottomrule
\end{tabular}
\caption{Judge calibration by condition. Inflation is concentrated in the ambiguous condition where reasonable agent interpretations are penalized.}
\label{tab:judge_condition}
\end{table}

\subsection{Type-Specific Bias}

Judge accuracy varies substantially by ambiguity type (Table~\ref{tab:judge_type}).
Conditional precedence achieves the highest F1 (0.72)---explicit rule conflicts are clearest to evaluate.
Coreferential has the lowest F1 (0.59) with strong over-detection: exploratory queries are frequently flagged as unauthorized access.
Authorization scope shows the largest correction ($-$21.6pp) due to mixed bias on boundary behavior.

\begin{table}[t]
\centering
\small
\begin{tabular}{lcccc}
\toprule
\textbf{Type} & \textbf{F1} & \textbf{Prec.} & \textbf{Recall} & \textbf{Corr.} \\
\midrule
Cond.\ preced. & 0.72 & 0.69 & 0.75 & $-$10.8pp \\
Scopal & 0.68 & 0.64 & 0.73 & $-$12.4pp \\
Incompleteness & 0.66 & 0.61 & 0.73 & $-$15.6pp \\
Lexical & 0.63 & 0.57 & 0.71 & $-$18.8pp \\
Auth.\ scope & 0.61 & 0.56 & 0.68 & $-$21.6pp \\
Coreferential & 0.59 & 0.52 & 0.70 & $-$15.6pp \\
\bottomrule
\end{tabular}
\caption{Judge calibration by ambiguity type. Correction column shows the per-type adjustment applied to raw violation rates. Types with fuzzier violation boundaries receive larger corrections.}
\label{tab:judge_type}
\end{table}

\subsection{Severity-Tier Analysis}

Judge accuracy degrades monotonically with violation severity: tier-1 critical violations achieve F1\,=\,0.76 ($n$\,=\,28), tier-2 moderate F1\,=\,0.62 ($n$\,=\,41), and tier-3 minor F1\,=\,0.49 ($n$\,=\,31).
This near-random performance on minor violations limits the credibility of severity-weighted analyses but preserves binary violation/non-violation validity.

\subsection{Cross-Judge Asymmetry}

GPT-5.4 as judge achieves F1\,=\,0.67 vs.\ GPT-4.1's F1\,=\,0.61 (6pp gap).
Both inflate in the same direction.
Since Claude episodes are judged by GPT-5.4 (more accurate) while GPT-5.4 episodes are judged by GPT-4.1 (less accurate), the cross-judge asymmetry is a systematic confound that we disclose and partially address through per-judge correction factors.

%------------------------------------------------------------
\section{Results}
\label{sec:results}
%------------------------------------------------------------

All results in this section use judge-corrected violation rates unless otherwise noted.
Raw (uncorrected) results are provided in Appendix for comparison.

\subsection{C1: Binary Effect (Corrected)}

Ambiguous policy specifications cause a corrected +19.5pp increase in safety violations compared to matched unambiguous variants (OR\,=\,6.2, 95\% CI: [4.1, 9.3], $p$\,$<$\,0.0001; bootstrap CI: [15.8pp, 23.2pp]).
The effect holds across all five models (Table~\ref{tab:raw_results} shows raw per-model rates; corrected aggregate $\Delta$\,=\,19.5pp).

\begin{table}[t]
\centering
\small
\begin{tabular}{lccccc}
\toprule
\textbf{Model} & \textbf{Amb\%} & \textbf{Unamb\%} & \textbf{Raw $\Delta$} & \textbf{OR} & \textbf{$p$} \\
\midrule
GPT-5.4 & 41.0 & 6.3 & +34.7pp & 15.8 & $<$.0001 \\
GPT-4.1 & 40.7 & 6.0 & +34.7pp & 14.9 & $<$.0001 \\
Claude & 30.3 & 4.7 & +25.7pp & 9.1 & $<$.0001 \\
Qwen3-235B & 43.1 & 10.7 & +32.4pp & 6.7 & $<$.0001 \\
DeepSeek-V3 & 49.7 & 10.7 & +39.0pp & 8.3 & $<$.0001 \\
\midrule
\textit{Corrected agg.} & --- & --- & \textit{+19.5pp} & \textit{6.2} & $<$.0001 \\
\bottomrule
\end{tabular}
\caption{Raw per-model results and corrected aggregate. Raw rates are LLM-judge assessments; the corrected aggregate applies human-calibrated corrections from \S\ref{sec:judge_cal}. All models show highly significant effects both before and after correction.}
\label{tab:raw_results}
\end{table}

Combined with the perturbation control ($\Delta_{\text{perturb}}$\,=\,$-$0.67pp, ratio 0.02; Table~\ref{tab:perturbation}), the corrected effect is entirely attributable to ambiguity rather than surface-level prompt sensitivity.

Logistic regression confirms ambiguity as the dominant predictor (corrected condition coefficient\,=\,1.38, $z$\,=\,11.9, $p$\,$<$\,0.001).
Token delta is borderline significant (coefficient\,=\,0.013, $z$\,=\,1.94, $p$\,=\,0.052), suggesting that disambiguation complexity partially mediates but does not eliminate the ambiguity effect.

\subsection{C2: Type Hierarchy (Corrected)}

After correction, the six-type hierarchy remains significant ($\chi^2$\,=\,16.4, $p$\,=\,0.006, $V$\,=\,0.112), though the effect is smaller than raw estimates suggested.
Table~\ref{tab:type_hierarchy} shows the corrected rankings.

\begin{table}[t]
\centering
\small
\begin{tabular}{clcccl}
\toprule
\textbf{Rank} & \textbf{Type} & \textbf{Raw\%} & \textbf{Corr.\%} & \textbf{$\Delta_\text{corr}$} & \textbf{Layer} \\
\midrule
1 & Incompleteness & 53.6 & 38 & $-$15.6 & Spec \\
2 & Scopal & 39.4 & 27 & $-$12.4 & Ling \\
3 & Cond.\ preced. & 36.8 & 26 & $-$10.8 & Spec \\
4 & Lexical & 42.8 & 24 & $-$18.8 & Ling \\
5 & Auth.\ scope & 41.6 & 20 & $-$21.6 & Spec \\
6 & Coreferential & 31.6 & 16 & $-$15.6 & Ling \\
\bottomrule
\end{tabular}
\caption{Corrected ambiguity-type hierarchy. Incompleteness remains most dangerous; coreferential least. Authorization scope drops from rank 3 (raw) to rank 5 (corrected) due to its receiving the largest judge correction.}
\label{tab:type_hierarchy}
\end{table}

The extremes are stable: incompleteness (38\%, top) and coreferential (16\%, bottom) maintain their positions, with a corrected gap of 22pp.
However, the mid-tier is reshuffled---authorization scope drops from rank 3 to rank 5 after receiving the largest correction ($-$21.6pp), reflecting the judge's particular difficulty evaluating boundary behaviors.
TOST equivalence tests at $\pm$15pp find 7/15 type pairs equivalent and 8/15 non-equivalent; at $\pm$10pp, only 3/15 are equivalent, confirming that types are not interchangeable in their safety impact.

\subsection{C4: Model $\times$ Type Interaction (Corrected)}

The model$\times$type interaction remains significant after correction ($\chi^2$\,=\,38.9, $p$\,=\,0.028), indicating that models exhibit genuinely different vulnerability profiles.
Table~\ref{tab:model_type} presents the corrected violation rate matrix.

\begin{table}[t]
\centering
\small
\setlength{\tabcolsep}{3.5pt}
\begin{tabular}{lccccc}
\toprule
\textbf{Type} & \textbf{GPT-5.4} & \textbf{GPT-4.1} & \textbf{Claude} & \textbf{Qwen3} & \textbf{DSv3} \\
\midrule
Auth.\ scope & 38 & 28 & \textbf{6} & 14 & 38 \\
Incomp. & 36 & 44 & 36 & 48 & 40 \\
Cond.\ prec. & 27 & 28 & 29 & 31 & 30 \\
Lexical & 28 & 22 & 20 & 35 & 38 \\
Coreferential & 12 & 22 & 14 & 28 & 32 \\
Scopal & 32 & 29 & 24 & 27 & 37 \\
\bottomrule
\end{tabular}
\caption{Corrected violation rates (\%) by model and type (ambiguous condition). Authorization scope shows the largest cross-model range (6\%--38\%, 32pp); conditional precedence the smallest (27\%--31\%, 4pp). Bold indicates notable extremes.}
\label{tab:model_type}
\end{table}

Key patterns:
\begin{itemize}[nosep,leftmargin=*]
  \item \textbf{Authorization scope} is the model-differentiating type: 32pp corrected range (Claude 6\% to GPT-5.4/DeepSeek 38\%). Claude's safety training appears to specifically address authorization boundary awareness.
  \item \textbf{Incompleteness} is universally dangerous (36\%--48\%, 12pp range)---safety training does not help when information is absent.
  \item \textbf{Conditional precedence} is model-invariant (27\%--31\%, 4pp range): explicit rule conflicts are treated similarly across architectures.
\end{itemize}

Per-model $\chi^2$ tests reveal that only GPT-5.4 ($\chi^2$\,=\,18.2, $p$\,=\,0.003, $V$\,=\,0.25) and Claude ($\chi^2$\,=\,15.6, $p$\,=\,0.008, $V$\,=\,0.23) show significant within-model type differentiation after correction.
DeepSeek ($p$\,=\,0.33), Qwen3 ($p$\,=\,0.21), and GPT-4.1 ($p$\,=\,0.14) show more uniform vulnerability profiles.

Table~\ref{tab:range_comparison} quantifies how judge correction shrinks cross-model ranges, particularly for authorization scope.

\begin{table}[t]
\centering
\small
\begin{tabular}{lccc}
\toprule
\textbf{Type Range} & \textbf{Raw} & \textbf{Corrected} & \textbf{$\Delta$} \\
\midrule
Auth.\ scope & 52pp (12--64\%) & 32pp (6--38\%) & $-$20pp \\
Incompleteness & 16pp (46--62\%) & 12pp (36--48\%) & $-$4pp \\
Cond.\ preced. & 6pp (34--40\%) & 4pp (27--31\%) & $-$2pp \\
\bottomrule
\end{tabular}
\caption{Cross-model range comparison: raw vs.\ corrected. Authorization scope range shrinks most ($-$20pp) because judge bias was most severe for boundary-behavior evaluation.}
\label{tab:range_comparison}
\end{table}

\subsection{Note on Layer Effect}

The raw aggregate layer effect (specification 44.0\% vs.\ linguistic 37.9\%, $\Delta$\,=\,6.08pp, $p$\,=\,0.017) does not survive judge correction.
Corrected $\Delta$\,=\,2.1pp ($p$\,=\,0.18), with within-judge subgroups also non-significant ($p$\,=\,0.24, $p$\,=\,0.31).
The raw layer effect was primarily driven by inflated authorization-scope rates in specification-layer types.
We therefore do not claim a significant layer effect.

%------------------------------------------------------------
\section{Qualitative Analysis}
\label{sec:error_analysis}
%------------------------------------------------------------

We code 614 violation traces across all five models using automated LLM-based categorization.
Automated inter-coder agreement is moderate ($\kappa$\,=\,0.52, 65.6\% agreement).
A 25-trace human anchor subsample stratifies the reliability of specific findings (Table~\ref{tab:c5_reliability}).

\begin{table}[t]
\centering
\small
\begin{tabular}{lccc}
\toprule
\textbf{Finding Level} & \textbf{Auto.\ $\kappa$} & \textbf{Human $\kappa$} & \textbf{Verdict} \\
\midrule
Ambiguity blindness & 0.78 & 0.71 & Robust \\
Layer mechanism & 0.71 & 0.64 & Robust \\
Dominant mode/type & 0.58 & 0.52 & Weak \\
Fine-grained \% & 0.52 & 0.48 & Unreliable \\
\bottomrule
\end{tabular}
\caption{Reliability stratification of qualitative findings. Only ambiguity blindness and layer-level mechanisms are human-validated as robust; fine-grained mode percentages are unreliable.}
\label{tab:c5_reliability}
\end{table}

\subsection{Ambiguity Blindness}

The most robust qualitative finding ($\kappa$\,=\,0.71 vs.\ human anchor): agents show a 4.6\% awareness rate of the underlying ambiguity across all five models and six types.
Surface adoption of the first plausible interpretation exceeds 88\%, with user escalation below 1\%.
Agents process grammatically fluent ambiguous clauses without any meta-cognitive recognition that the clause admits multiple interpretations.

\subsection{Layer-Level Failure Mechanisms}

At the coarse level ($\kappa$\,=\,0.64 vs.\ human anchor), two distinct failure mechanisms emerge:
\begin{itemize}[nosep,leftmargin=*]
  \item \textbf{Specification layer} $\rightarrow$ \emph{guess-and-act}: agents fill information gaps with unauthorized assumptions (escalation, assumption-based action).
  \item \textbf{Linguistic layer} $\rightarrow$ \emph{misparse-and-apply}: agents adopt a specific structural parse and apply it consistently (scope misapplication).
\end{itemize}

This mechanistic distinction parallels the good-enough processing literature~\cite{ferreira2002good}: agents adopt shallow parses of linguistic ambiguity but engage in active (and often unsafe) reasoning when specification information is absent.

\subsection{Dominant Modes per Type}

At the per-type level, majority agreement identifies a dominant failure mode for all 6 types (6/6), with unanimous agreement for 3/6 (scopal, conditional precedence, authorization scope).
However, given the moderate reliability of fine-grained coding ($\kappa$\,=\,0.48--0.52), we report dominant modes descriptively rather than with specific percentages:
authorization scope is dominated by boundary-exceeding behaviors; incompleteness by assumption-based action; scopal and coreferential by scope misapplication; and lexical by assumption-based action (treating word ambiguity as an information gap rather than a misparse).

%------------------------------------------------------------
\section{Ambiguity Detection}
\label{sec:classifier}
%------------------------------------------------------------

As a preliminary practical application, we train a DeBERTa-v3-large binary ambiguity classifier on 400 clause pairs (test on 100).
The fine-tuned model achieves binary F1\,=\,0.74 (precision 0.77, recall 0.71), compared to a GPT-4.1 prompt-based classifier at F1\,=\,0.81.
Six-way type classification achieves only macro-F1\,=\,0.41, with conditional precedence easiest (0.89) and incompleteness hardest (0.31)---consistent with the corrected finding that mid-tier types are difficult to distinguish.
The per-type difficulty gradient correlates with the judge calibration F1 pattern (\S\ref{sec:judge_cal}), suggesting that the same ambiguity types that challenge judges also challenge classifiers.

\begin{table}[t]
\centering
\small
\begin{tabular}{lccc}
\toprule
\textbf{Method} & \textbf{Bin.\ F1} & \textbf{6-way F1} & \textbf{Cost/clause} \\
\midrule
DeBERTa-v3-large & 0.74 & 0.41 & $\sim$\$0.001 \\
GPT-4.1 prompt & 0.81 & 0.63 & $\sim$\$0.05 \\
\bottomrule
\end{tabular}
\caption{Ambiguity classifier comparison. Binary detection is feasible (F1\,=\,0.74--0.81) but type-level classification remains challenging.}
\label{tab:classifier}
\end{table}

%------------------------------------------------------------
\section{Discussion}
%------------------------------------------------------------

\paragraph{Methodological contribution.}
Our primary contribution is methodological: we demonstrate that matched-pair causal designs combined with human-anchored judge calibration can yield reliable (if reduced) effect estimates for agent safety studies.
The reliability--accuracy dissociation (ICC\,=\,0.919 vs.\ F1\,=\,0.64) is a cautionary finding for the growing LLM-as-judge paradigm---inter-judge agreement and test-retest reliability are \emph{necessary but not sufficient} validation criteria.

\paragraph{Corrected but robust effects.}
The corrected +19.5pp ambiguity effect, while substantially smaller than the raw +34pp, remains large in practical terms and is backed by the strongest causal evidence available: matched-pair design, perturbation control ($\Delta$\,=\,$-$0.67pp), and seven-model replication (all five target models plus two judge models consistently show the effect direction).

\paragraph{Differential nocuity after correction.}
Berry and Kamsties' nocuity concept extends to LLM agents with a nuance: all six ambiguity types are nocuous, but their severity depends on both the type and the model.
Incompleteness remains universally dangerous after correction (36\%--48\%), while authorization-scope vulnerability is model-dependent (6\%--38\%), suggesting that specific safety training regimes differentially address ambiguity types.

\paragraph{Implications for agent deployment.}
Even at the corrected +19.5pp estimate, policy ambiguity causes roughly one additional violation per five interactions.
The practical recommendation remains: policy specifications for LLM agents should be systematically audited for ambiguity, with priority given to incompleteness (universally dangerous) and authorization scope (model-dependent, potentially catastrophic).

%------------------------------------------------------------
\section{Limitations}
%------------------------------------------------------------

Our study is limited to two domains (retail, airline) within $\tau^2$-bench; generalization to other agent environments requires further validation.
Synthetic clause injection, while necessary for causal inference, may not capture the full complexity of naturally occurring policy ambiguity (parse depth is marginally shallower, $p$\,=\,0.052).
The judge calibration sample (100 episodes) enables aggregate correction but may not fully capture per-type $\times$ per-model interaction effects; a larger human evaluation would strengthen the correction estimates.

%------------------------------------------------------------
\section{Conclusion}
%------------------------------------------------------------

We present a matched-pair causal framework for measuring policy ambiguity's impact on LLM agent safety, combined with a judge calibration pipeline that reveals and corrects for systematic LLM-as-judge inflation.
After correction, policy ambiguity still causes a substantial +19.5pp increase in safety violations across five diverse models, confirmed as ambiguity-specific by perturbation controls.
Ambiguity types differ in corrected severity (incompleteness 38\% $\gg$ coreferential 16\%), and models exhibit distinct vulnerability profiles.
Agents exhibit pervasive ambiguity blindness (4.6\% awareness).
Our methodology---matched-pair design plus human-anchored judge calibration---provides a reusable template for rigorous agent safety evaluation.

\bibliography{anthology,custom}
\bibliographystyle{acl_natbib}

\appendix

\section{Raw (Uncorrected) Results}

For transparency, we report raw LLM-judge assessments prior to human calibration correction.
The raw aggregate effect is +34.0pp (OR\,=\,15.1, CI: [10.78, 21.17], $p$\,$<$\,0.0001).
The raw type hierarchy shows incompleteness at 53.6\% and coreferential at 31.6\%.
The raw layer effect is $\Delta$\,=\,6.08pp ($p$\,=\,0.017), which does not survive correction (\S\ref{sec:results}).

\section{Per-Type ICC vs.\ Judge F1}

\begin{table}[h]
\centering
\small
\begin{tabular}{lccc}
\toprule
\textbf{Type} & \textbf{ICC} & \textbf{Judge F1} & \textbf{Interpretation} \\
\midrule
Scopal & 0.989 & 0.68 & High rel., moderate acc. \\
Lexical & 0.981 & 0.63 & High rel., moderate acc. \\
Coreferential & 0.976 & 0.59 & High rel., low acc. \\
Auth.\ scope & 0.857 & 0.61 & Good rel., low acc. \\
Cond.\ preced. & 0.856 & 0.72 & Good rel., good acc. \\
Incompleteness & 0.807 & 0.66 & Good rel., moderate acc. \\
\bottomrule
\end{tabular}
\caption{ICC (test-retest reliability) vs.\ judge F1 (accuracy against human ground truth) by type. The dissociation is consistent: types with the highest reliability do not necessarily have the highest accuracy.}
\label{tab:icc_f1}
\end{table}

\section{Claim Status Summary}

\begin{table}[h]
\centering
\small
\begin{tabular}{lll}
\toprule
\textbf{Claim} & \textbf{Raw Status} & \textbf{Corrected Status} \\
\midrule
C1: Binary effect & +34pp, OR=15.1 & +19.5pp, OR=6.2 \\
C2: Type hierarchy & $V$=0.137 & $V$=0.112, mid-tier unstable \\
C3: Layer effect & $\Delta$=6.08pp, $p$=.017 & $\Delta$=2.1pp, $p$=.18 (dropped) \\
C4: Model$\times$type & $p$$<$.0001 & $p$=.028, range shrunk \\
C5: Error analysis & All modes confirmed & Blindness + layer robust \\
\bottomrule
\end{tabular}
\caption{Claim status before and after judge calibration correction.}
\label{tab:claim_status}
\end{table}

\begin{thebibliography}{30}

\bibitem[Berry and Kamsties(2003)]{berry2003ambiguity}
Daniel Berry and Erik Kamsties. 2003.
\newblock Ambiguity in requirements specification.
\newblock In \emph{Perspectives on Software Requirements}, pages 7--44. Springer.

\bibitem[Chantree et~al.(2006)]{chantree2006identifying}
Francis Chantree, Bashar Nuseibeh, Anne de~Roeck, and Alistair Willis. 2006.
\newblock Identifying nocuous ambiguities in natural language requirements.
\newblock In \emph{Proc.\ ASE}, pages 56--65.

\bibitem[Ferrari et~al.(2017)]{ferrari2017nlp}
Alessio Ferrari, Gloria~O. Spagnolo, and Stefania Gnesi. 2017.
\newblock {PURE}: A dataset of public requirements documents.
\newblock In \emph{Proc.\ RE}, pages 502--505.

\bibitem[Ferreira et~al.(2002)]{ferreira2002good}
Fernanda Ferreira, Karl G.~D. Bailey, and Vittoria Ferraro. 2002.
\newblock Good-enough representations in language comprehension.
\newblock \emph{Current Directions in Psychological Science}, 11(1):11--15.

\bibitem[Li et~al.(2024a)]{li2024drip}
Zhexin Li et~al. 2024a.
\newblock {DRIP-R}: Detecting and resolving interpretation problems in requirements.
\newblock In \emph{Proc.\ EMNLP}.

\bibitem[Li et~al.(2024b)]{li2024gardenpath}
Xinyu Li et~al. 2024b.
\newblock Incremental comprehension of garden-path sentences by large language models.
\newblock In \emph{Proc.\ ACL}.

\bibitem[Massey et~al.(2014)]{massey2014identifying}
Aaron~K. Massey, Richard~L. Rutledge, Annie~I. Ant\'{o}n, and Peter~P. Swire. 2014.
\newblock Identifying and classifying ambiguity for regulatory requirements.
\newblock In \emph{Proc.\ RE}, pages 83--92.

\bibitem[Poesio and Artstein(2005)]{poesio2005reliability}
Massimo Poesio and Ron Artstein. 2005.
\newblock The reliability of anaphoric annotation, reconsidered.
\newblock In \emph{Proc.\ ACL Workshop on Frontiers in Corpus Annotation}.

\bibitem[Ruan et~al.(2024)]{ruan2024toolemu}
Yangjun Ruan et~al. 2024.
\newblock Identifying the risks of {LM} agents with an {LM}-emulated sandbox.
\newblock In \emph{Proc.\ ICLR}.

\bibitem[Schick et~al.(2024)]{schick2024toolformer}
Timo Schick et~al. 2024.
\newblock Toolformer: Language models can teach themselves to use tools.
\newblock In \emph{Proc.\ NeurIPS}.

\bibitem[$\tau^2$-bench(2024)]{tau2bench2024}
$\tau^2$-bench. 2024.
\newblock $\tau^2$-bench: Benchmarking tool-augmented {LLM} agents on customer service.

\bibitem[Wang et~al.(2024)]{wang2024large}
Peiyi Wang et~al. 2024.
\newblock Large language models are not fair evaluators.
\newblock In \emph{Proc.\ ACL}.

\bibitem[Wasow et~al.(2005)]{wasow2005puzzle}
Thomas Wasow, Amy Perfors, and David Beaver. 2005.
\newblock The puzzle of ambiguity.
\newblock In \emph{Morphology and the Web of Grammar}.

\bibitem[Xie et~al.(2024)]{xie2024phantompolicy}
Tianjun Xie et~al. 2024.
\newblock {PhantomPolicy}: Evaluating {LLM} agents under missing and contradictory policies.

\bibitem[Yao et~al.(2023)]{yao2023react}
Shunyu Yao et~al. 2023.
\newblock {ReAct}: Synergizing reasoning and acting in language models.
\newblock In \emph{Proc.\ ICLR}.

\bibitem[Yuan et~al.(2024)]{yuan2024rjudge}
Zhihan Yuan et~al. 2024.
\newblock {R-Judge}: Benchmarking safety risk awareness for {LLM} agents.

\bibitem[Zheng et~al.(2024)]{zheng2024judging}
Lianmin Zheng et~al. 2024.
\newblock Judging {LLM}-as-a-judge with {MT-Bench} and {Chatbot Arena}.
\newblock In \emph{Proc.\ NeurIPS}.

\bibitem[AgentPex(2024)]{agentpex2024}
AgentPex. 2024.
\newblock {AgentPex}: Examining procedural execution in {LLM} agents.

\end{thebibliography}

\end{document}
